\ifdefined\iscombined
	\sectiontitle{Section 1 - Practice Problems}
	\setcounter{section}{1}
  % being included -> skip preamble
\else
\documentclass[10pt]{article}
\usepackage{tabularx}
\usepackage{mathtools}
\usepackage{amsmath, amssymb, amsthm}
\usepackage{enumitem}
\usepackage{geometry}
\usepackage{fancyhdr}
\usepackage{titlesec}
\usepackage{hyperref}
\geometry{margin=1in}

\pagestyle{fancy}
\fancyhf{}
\title{Section 1 - Practice Problems}
\author{Matthew Farah, \href{mailto:farahm11@mcmaster.ca}{$\langle$farahm11@mcmaster.ca$\rangle$}, 400468251}
\date{\today}

\hypersetup{
        colorlinks=true,
        linkcolor=blue,
        filecolor=magenta,
        urlcolor=blue,
        citecolor=blue,
	anchorcolor=blue
}

\fancyhead{}
\fancyhead[L]{Matthew Farah, \href{mailto:farahm11@mcmaster.ca}{$\langle$farahm11@mcmaster.ca$\rangle$}, 400468251}
\fancyhead[R]{Section 1 - Practice Problems}

\newcommand{\chaptertitle}[1]{
	\refstepcounter{chapter}
	\vspace{1.5em}
	\begin{center}
	\underline{\noindent\Large\bfseries{Chapter} \thechapter: #1}
	\end{center}
	\vspace{1em}
}

\newcommand{\sectiontitle}{Section 1 - Practice Problems}
\setcounter{section}{1}

\newcounter{chapter}[section]
\renewcommand{\thechapter}{\thesection.\arabic{chapter}}

\newcounter{exercise}[chapter]
\renewcommand{\theexercise}{\thechapter.\arabic{exercise}}

\newenvironment{exercise}[1][]{
	\ifx\\#1\\
		\refstepcounter{exercise}
	\else
		\setcounter{exercise}{#1}
	\fi
	\large{\par\noindent\textbf{Exercise \theexercise}}
}{\vspace{0.8em}}

\newenvironment{solution}{
	\vspace{0.4cm}\par\noindent\large\textbf{{Solution:}}\par\vspace{0.4em}
}{\vspace{0.8em}}

\newcounter{partslist}

\makeatletter
\newcommand{\parts@seriesname}{parts@\theexercise}

\newenvironment{parts}{
	\edef\parts@ser{\parts@seriesname}
	\ifcsname\parts@ser\endcsname
		\begin{enumerate}[resume=\parts@ser,label=\textbf{(\alph*)},leftmargin=3.3em,itemsep=0.4em,before=\large]
	\else
		\begin{enumerate}[series=\parts@ser,label=\textbf{(\alph*)},leftmargin=3.3em,itemsep=0.4em,before=\large]
		\expandafter\gdef\csname\parts@ser\endcsname{1}
	\fi
}{
	\end{enumerate}
}
\makeatother

\makeatletter
\newcommand{\steps@seriesname}{steps@\theexercise}

\newenvironment{steps}{
	\edef\steps@ser{\steps@seriesname}
	\ifcsname\steps@ser\endcsname
		\begin{enumerate}[resume=\steps@ser,label=(\Roman*),leftmargin=3.15em,itemsep=0.4em,topsep=0.6em]
	\else
		\begin{enumerate}[series=\steps@ser,label=(\Roman*),leftmargin=3.15em,itemsep=0.4em,topsep=0.6em]
		\expandafter\gdef\csname\steps@ser\endcsname{1}
	\fi
}{
	\end{enumerate}
}

\makeatother

\newcolumntype{Y}{>{\centering\arraybackslash}X}
\renewcommand{\arraystretch}{1.8}

\fontsize{10pt}{14pt}\selectfont

\begin{document}

\maketitle

\fi

\vspace{-1cm}


\chaptertitle{Sets and Functions}

\begin{exercise}[5]
	Prove the following distributive laws:
\end{exercise}

\begin{parts}
	\item $A \cap (B \cup C) = (A \cap B) \cup (A \cap C)$.
\end{parts}

\begin{solution}
	We can prove that two sets are equal by showing that they are both subsets of one another.

	\vspace{.25cm}

	\begin{steps}
		\item Show that $A \cap (B \cup C) \subseteq (A \cap B) \cup (A \cap C)$.
	\end{steps}

	\vspace{.25cm}

	Let $x \in (A \cap (B \cup C))$. Therefore, by definition of intersection and union, we know $x \in A \; \wedge \; x \in (B \cup C)$. Thus, this implies that $x \in B \vee x \in C$. However, $x \in A \implies x \in ((A \cap B) \cup (A \cap C)) \iff x \in (B \cup C)$. However, we already determined that $x \in B \vee x \in C$. Thus, $x \in (A \cap (B \cup C)) \implies x \in ((A \cap B) \cup (A \cap C))$ and therefore $A \cap (B \cup C) \subseteq (A \cap B) \cup (A \cap C)$.

	\vspace{.25cm}

	\begin{steps}
		\item Show that $A \cap (B \cup C) \supseteq (A \cap B) \cup (A \cap C)$.
	\end{steps}

	\vspace{.25cm}

	Let $x \in ((A \cap B) \cup (A \cap C))$. Therefore, by definition of intersection and union, we know $(x \in A \wedge x \in B) \wedge (x \in A \wedge x \in C)$. This clearly implies that $x \in A$, and furthermore that $x \in C \vee x \in B$. Therefore, we know $x \in (B \cup C)$, and since $x \in A$, $x \in A \cap (B \cup C)$. Thus, $x \in ((A \cap B) \cup (A \cap C)) \implies x \in (A \cap (B \cup C))$ and thus $A \cap (B \cup C) \supseteq (A \cap B) \cup (A \cap C)$.

	\vspace{.75cm}

	Therefore, since $A \cap (B \cup C) \subseteq (A \cap B) \cup (A \cap C)$ and $A \cap (B \cup C) \supseteq (A \cap B) \cup (A \cap C)$, $A \cap (B \cup C) = (A \cap B) \cup (A \cap C)$.
\end{solution}

\begin{parts}
	\item $A \cup (B \cap C) = (A \cup B) \cap (A \cup C)$.
\end{parts}

\begin{solution}
	We will prove this in much the same way as in \textbf{(a)}.

	\begin{steps}
		\item Show that $A \cup (B \cap C) \subseteq (A \cup B) \cap (A \cup C)$.
	\end{steps}

	Let $x \in (A \cup (B \cap C))$. Therefore, $x \in A \vee x \in (B \cap C)$. If $x \in (B \cap C)$ which means we know that $x \in B \wedge x \in C$. Therefore, we know that $x \in A \vee (x \in B \wedge x \in C)$. If $x \in A$, then $x \in (A \cup B) \wedge x \in (A \cup C)$, and therefore $x \in ((A \cup B) \cap (A \cup C))$, and if $x \in B \wedge x \in C$, then $x \in (A \cup B) \wedge x \in (A \cup C)$ and therefore $x \in ((A \cup B) \cap (A \cup C))$. Therefore, $x \in (A \cup (B \cap C)) \implies x \in ((A \cup B) \cap (A \cup B)$ and thus $A \cup (B \cap C) \subseteq ((A \cup B) \cap (A \cup C))$.

	\begin{steps}
		\item Show that $A \cup (B \cap C) \supseteq (A \cup B) \cap (A \cup C)$.
	\end{steps}

	Let $x \in ((A \cup B) \cap (A \cup C))$. Therefore, $x \in (A \cup B) \wedge x \in (A \cup C)$ which means we know that $x \in A \vee (x \in B \wedge x \in C)$. If $x \in A$, then trivially $x \in (A \cup (B \cap C))$, and if $x \in B \wedge x \in C$, then $x \in (B \cap C)$ and therefore $x \in (A \cup (B \cap C))$. Therefore $x \in ((A \cup B) \cap (A \cup B)) \implies x \in (A \cup (B \cap C))$ and thus $(A \cup B) \cap (A \cup C) \subseteq A \cup (B \cap C)$.

	\vspace{.75cm}

	Therefore, since $A \cup (B \cap C) \subseteq (A \cup B) \cap (A \cup C)$ and $A \cup (B \cap C) \supseteq (A \cup B) \cap (A \cup C)$, $A \cap (B \cap C) = (A \cup B) \cap (A \cup C)$.
\end{solution}


\begin{exercise}[11]
	Let $g(x) \coloneq x^2$ and $f(x) \coloneq x + 2$ for $x \in \mathbb{R}$, and let $h$ be the composite function $h \coloneq g \circ f$.
\end{exercise}

\begin{parts}
	\item Find the direct image $h(E)$ of $E \coloneq \{x \in \mathbb{R}: 0 \le x \le 1\}$.
\end{parts}

\begin{solution}
	Recall that the image of a function under a set is the set of all values which the function can attain with inputs taken from the set. We know the composite function $h$ is given by $h(x) = {f(x)}^2 = (x + 2)^2$. Therefore, the image of $h$ under $E$ is $\{y \in \mathbb{R}: 4 \le y \le 9 \}$.
\end{solution}

\begin{parts}
	\item Find the inverse image $h^{-1}(G)$ of $G \coloneq \{x \in \mathbb{R}: 0 \le x \le 4\}$.
\end{parts}

\begin{solution}
	Recall that the inverse image of a function is the set of all values in the domain of the function which are in the given set. Since we previously determined that $h(x) = (x + 2)^2$, by rearranging, we can find that:

	\begin{align*}
		x &= (h^{-1}(x) + 2)^2 \\
		\pm \sqrt{x} &= (h^{-1}(x) + 2) \\
		-2 \pm \sqrt{x} &= h^{-1}(x) \\
	\end{align*}

	Therefore, $h{-1}(x) = -2 \pm \sqrt{x}$. Thus, the inverse image of $h$ under $G$ is $\{x \in \mathbb{R}: -2 \le x \le 2 \}$.
\end{solution}

\begin{exercise}[19]

\end{exercise}

\begin{parts}
	\item Show that if $f: A \rightarrow B$ is injective and $E \subseteq A$, then $f^{-1}(f(E)) = E$. Give an example to show that equality need not hold if $f$
	is not injective.
	
	\begin{solution}
		\begin{proof}
			Let $x \in E$. By definition of inverse, we know trivially that $x \in E \implies x \in f^{-1}(f(E))$ and thus $E \subseteq f^{-1}f(E)$. Now consider $x \in f^{-1}(f(E))$. To prove that $x$ must necessarily be in $E$, we rely on a proof by contradiction. Assume for the purposes of contradiction that $x \notin f^{-1}(f(E))$. Since we are taking the inverse of $f(E)$, that implies $\exists y \in E$ such that $f(y) = f(x)$. But since $x \notin E$ and $y \in E$, we necessarily know that $x \neq y$. However, this would violate the assumption that $f$ is injective since, if this were true $\exists x, y \in A, x \neq y$ such that $f(x) = f(y)$. Thus, by contradiction, $x \in f^{-1}(f(E)) \implies x \in E$ and $f^{-1}(f(E)) \subseteq E$.

			Therefore, since $E \subseteq f^{-1}(f(E))$ and $f^{-1}(f(E)) = \subseteq E$, $E = f^{-1}(f(x))$ for injective $f$.
		\end{proof}

		If injectivity is not assumed, this is not necessarily true. Taking $f(x) = x^2$ and $E = [0, 1]$ as an example, $f^{-1}(f(x)) = [-1, 1]$. Clearly then, for non-injective $f$, equality need not hold.
	\end{solution}
\end{parts}

\begin{parts}
	\item Show that if $f: A \rightarrow B$ is surjective and $H \subseteq B$, then $f(f^{-1}(H)) = H$. Give an example to show that equality need not hold if $f$ is not surjective.
\end{parts}

\begin{solution}
	
	\begin{proof}
		Let $K = f^{-1}(H)$. Since $f$ is surjective, we know that $f(K) = H$ and since $K = f^{-1}(H)$, then $H = f(f^{-1}(H))$. Thus, if $x \in f(f^{-1}(H))$, then $x \in H$ and therefore $H \subseteq f(f^{-1}(H))$. \\
	\end{proof}
	
	Kind of a gross counter example but i'm tired sorry.

	If surjectivity is not assumed, this is not necessarily true. Once again, take $f(x) = x^2$  and $H = (-\infty, \infty)$. $f(f^{-1}(H)) = [0, \infty)$. Clearly then, for non-surjective $f$, equaltiy need not hold.
\end{solution}

\begin{exercise}[21]
	Prove that if $f: A \rightarrow B$ and $g: B \rightarrow C$ are bijective, then the composite $g \circ f$ is a bijective map of $A$ onto $C$.
\end{exercise}

\begin{solution}
	\begin{proof}
		To prove that $g \circ f$ is a bijective map of A onto C, we must show that $g \circ f$ is both injective and surjective.

		\begin{steps}
			\item Prove $g \circ f$ is injective.
		\end{steps}

		Since it is given that $f$ and $g$ are injective, then we know $\forall i, j \in A, i \neq j$, $f(i) \neq f(j)$ and $\forall k, l \in B$, $g(k) \neq g(l)$. Thus:

		\begin{align*}
			\forall x, y \in A, x \neq y &\implies \\
			f(x) \neq f(y) &\implies \\
			g(f(x)) \neq g(f(y))&
		\end{align*}

		Therefore, $g \circ f$ is an injective map of $A$ onto $C$.

		\begin{steps}
			\item Prove $g \circ f$ is surjective.
		\end{steps}

		Since it is given that $f$ and $g$ are surjective, $i \in B \implies \exists \, j \in A$ such that $f(j) = i$ and $k \in C \implies \exists \, l \in B$ such that $g(l) = k$. Thus:

		\begin{align*}
			x &\in C \implies \\
			\exists \; y \in B \ \text{such that} \ g(y) &= x \implies \\
			\exists \; z \in A \ \text{such that} \ f(z) &= x \implies \\
			g(f(z)) &= x
		\end{align*}

		Therefore $g \circ f$ is a surjective map of $A$ onto $C$.

		Thus, since $g \circ f$ is both an injective and surjective map of $A$ onto $C$, $g \circ f$ is bijective.

	\end{proof}
\end{solution}

\newpage

\chaptertitle{Mathematical Induction}

\begin{exercise}[1]
	Prove that $\frac{1}{1 \cdot 2} + \frac{1}{2 \cdot 3} + \ldots \frac{1}{n \cdot (n + 1)} = \frac{n}{n + 1} \quad \forall n \in \mathbb{N}$.
\end{exercise}

\begin{solution}
	
	\begin{proof}
		Let $P(n)$ be the statement $P(n): \frac{1}{1} \cdot 2 + \frac{1}{2} \cdot 3 + \ldots \frac{1}{n \cdot (n + 1)} = \frac{n}{n + 1}$.

		\begin{steps}
			\item Show the base case holds (Prove $P(1)$ is true).
		\end{steps}

		\begin{align*}
			&\text{LHS:}& &\text{RHS:} \\
			&\frac{1}{n \cdot (n + 1)}& &\frac{n}{n + 1} \\
			&=\frac{1}{1 \cdot (1 + 1)}& &=\frac{1}{1 + 1} \\
			&=\frac{1}{2}& &=\frac{1}{2}
		\end{align*}

		Thus, $P(1)$ is true.

		\begin{steps}
			\item State induction hypothesis (State $P(k)$).
		\end{steps}
		
		Assuming there exists a $k \in \mathbb{N}$ for which $P(k)$ is true we can formulate our induction hypothesis as:
		
		\begin{align*}
			P(k): \frac{1}{1} \cdot 2 + \frac{1}{2} \cdot 3 + \ldots \frac{1}{k \cdot (k + 1)} &= \frac{k}{k + 1}
		\end{align*}

		\newpage

		\begin{steps}
			\item Prove $P(k) \implies P(k + 1)$.
		\end{steps}

		\begin{align*}
			\frac{1}{1} \cdot 2 + \frac{1}{2} \cdot 3 + &\ldots \frac{1}{k \cdot (k + 1)} + \frac{1}{(k + 1) \cdot (k + 2)} \\\\
			&= \frac{k}{k + 1} + \frac{1}{(k + 1) \cdot (k + 2)} \\\\
			&= \frac{k(k + 2) + 1}{(k + 1) \cdot (k + 2)} \\\\
			&= \frac{k^2 + 2k + 1}{(k + 1) \cdot (k + 2)} \\\\
			&= \frac{(k + 1)^2}{(k + 1) \cdot (k + 2)} \\\\
			&= \frac{k + 1}{k + 2} \\
		\end{align*}

		Therefore, since $P(1)$ is true, and $P(k) \implies P(k + 1)$ for arbitrary $k \in \mathbb{N}$, $P(n)$ is true $\forall n \in \mathbb{N}$.
	\end{proof}
\end{solution}

\begin{exercise}[7]
	Prove that $5^{2n} - 1$ is divisible by $8 \; \forall n \in \mathbb{N}$.
\end{exercise}

\begin{solution}
	\begin{proof}
		Let $P(n)$ be the statement $P(n): 8 \; \text{mod} \; (5^{2n} - 1) = 0$.

		\begin{steps}
			\item Show the base case holds (Prove $P(1)$ true).
		\end{steps}

		\begin{align*}
			8 \; \text{mod} \; (5^{2n} - 1) &= \text{mod}(5^{2 \cdot 1} - 1) \\
			&= 8 \; \text{mod} \; (5^2 - 1) \\
			&= 8 \; \text{mod} \; (24) \\
			&= 0
		\end{align*}

		Therefore, $P(1)$ is true.

		\begin{steps}
			\item State induction hypothesis (State $P(k)$).
		\end{steps}

		Assuming there exists a $k \in \mathbb{N}$ for which $P(k)$ is true we can formulate our induction hypothesis as:

		\begin{align*}
			P(k): 8 \; \text{mod} \;(5^{2k} - 1)
		\end{align*}

		\begin{steps}
			\item Prove $P(k) \implies P(k + 1)$.
		\end{steps}

		\begin{align*}
			8 \; \text{mod} \;(5^{2(k + 1)} - 1) &= 8 \; \text{mod} \; (5^{2k + 2} - 1) \\
			&= 8\; \text{mod} \; (5^2 \cdot 5^2k - 1) \\
			&= 0 \\
		\end{align*}

		Therefore, since $P(1)$ is true, and $P(k) \implies P(k + 1)$ for arbitrary $k \in \mathbb{N}$, $P(n)$ is true $\forall n \in \mathbb{N}$.
	\end{proof}
\end{solution}

\begin{exercise}[11]
	Conjecture a formula for the sum of the first $n$ odd natural numbers $1 + 3 + \ldots + (2n - 1)$, and prove your formula using mathematical induction
\end{exercise}

\begin{solution}
	\begin{tabularx}{\linewidth}{ c || Y Y Y Y Y Y }
		First $n$ odd natural numbers & $1$ & $2$ & $3$ & $4$ & $5$ & $\ldots$ \\
		\hline
		$1 + 3 + \ldots + (2n - 1)$ & $1$ & $4$ & $9$ & $16$ & $25$ & $\ldots$ \\
	\end{tabularx}

	\vspace{.75cm}

	Based on this, I conjecture that $1 + 3 + \ldots + (2n - 1) = n^2$. Next, I will prove this fact.

	\begin{proof}
		Let $P(n)$ be the statement $P(n): 1 + 3 + \ldots + (2n - 1) = n^2$.

		\begin{steps}
			\item Show base case holds (Prove $P(1)$ true).
		\end{steps}

		\begin{align*}
			&\text{LHS:}& &\text{RHS:} \\
			&(2n - 1)& &n^2 \\
			&= (2(1) - 1)& &(1)^2 \\
			&=1& &=1
		\end{align*}

		Therefore, $P(1)$ is true.
		
		\begin{steps}
			\item State induction hypothesis (State $P(k)$).
		\end{steps}

		Assuming there exists a $k \in \mathbb{N}$ for which $P(k)$ is true we can formulate our induction hypothesis as:

		\begin{align*}
			P(k): 1 + 3 + \ldots + (2k - 1) = k^2
		\end{align*}

		\begin{steps}
			\item Prove $P(k) \implies P(k + 1)$
		\end{steps}

		\begin{align*}
			1 + 3 + \ldots + (2k - 1) + (2k + 1) &= k^2 + (2k + 1) \\
			&= (k + 1)^2 \\
		\end{align*}

		Therefore, since $P(1)$ is true, and $P(k) \implies P(k + 1)$ for arbitrary $k \in \mathbb{N}$, $P(n)$ is true $\forall n \in \mathbb{N}$.
	\end{proof}
\end{solution}

\begin{exercise}[13]
	Prove that $n < 2^n \; \forall n \in \mathbb{N}$.
\end{exercise}

\begin{solution}
	Let $P(n)$ be the statement $P(n): n < 2^n$.

	\begin{proof}

		\begin{steps}
			\item Show base case holds (Prove $P(1)$ true).
		\end{steps}

		\begin{align*}
			&\text{LHS:}& &\text{RHS:} \\
			&n& &2^n \\
			&=(1)& &=2^(1) \\
			&=1& &=2
		\end{align*}

		Therefore, $P(1)$ is true.

		\begin{steps}
			\item Formulate induction hypothesis (State $P(k)$).
		\end{steps}

		Assuming there exists a $k \in \mathbb{N}$ for which $P(k)$ is true we can formulate our induction hypothesis as:

		\begin{align*}
			P(k): k < 2^k
		\end{align*}

		\begin{steps}
			\item Prove $P(k) \implies P(k + 1)$
		\end{steps}

		\begin{align*}
			k + 1 &< 2^k + 1 \\
			&< 2^{k + 1} \\
		\end{align*}

		Therefore, since $P(1)$ is true, and $P(k) \implies P(k + 1)$ for arbitrary $k \in \mathbb{N}$, $P(n)$ is true $\forall n \in \mathbb{N}$.
	\end{proof}
\end{solution}

\newpage

\chaptertitle{Finite and Infinite Sets}

\begin{exercise}[3]
	Let $S \coloneq \{1, 2\}$ and $T \coloneq \{a, b,c \}$
\end{exercise}

\begin{parts}
	\item Determine the number of different injections from $S$ into $T$.
\end{parts}

\begin{solution}
	My answer is $6$ and my reasoning is as follows: \\

	Taking $s_1 \in S$, we have three possible choices of $t_1 \in T$ for which we can create an injection from $s_1$ to $T$, giving us 3 different injections. Now taking $s_2 \in S$, we would have three possible choices of $t_2 \in T$, but since the relation must be injective, $t_1$ cannot equal $t_2$ and thus we only have two possible choices. Thus, we have $3 \cdot 2 = 6$ possible injections from $S$ to $T$.
\end{solution}

\begin{parts}
	\item Determine the number of different surjections from $T$ onto $S$.
\end{parts}

\begin{solution}
	My answer is $6$ and my reasoning is as follows: \\

	Taking $t_1 \in T$, $t_1$ has two possible choices of $s_1 \in S$ for which we can create a surjection from $t_1$ to $S$. Now consider $t_2 \in T$. Since the relationship need not be injective, we separate our choices into two possible cases. If the surjection from $t_2$ to $S$ is $s_1$, then we know that $t_3$ \underline{MUST} be $s_2$. However, if the surjection from $t_2$ to $S$ is not $s_1$, then $t_3$ can relate to either of $s_1$ or $s_2$. Thus, we have $2 \cdot 1 \cdot 1 + 2 \cdot 1 \cdot 2 = 6$ possible surjections from $T$ to $S$.
\end{solution}

\begin{exercise}[5]
	Give an explicit definition of the bijection $f$ from $\mathbb{N}$ onto $\mathbb{Z}$ described in example 1.3.7.(b).
\end{exercise}

\begin{solution}
	\begin{align*}
		f(n) &= \begin{cases}
			0, \; \text{if} \; n = 1 \\
			\frac{n}{2}, \; \text{if} \; \text{n is even} \\
			-\frac{n - 1}{2}, \; \text{if} \; \text{n is odd and} \; n \neq 1
		\end{cases}
	\end{align*}
\end{solution}

\begin{exercise}[7]
	Prove that a set $T_1$ is denumerable if and only if there is a bijection from $T_1$ onto a denumerable set $T_2$.
\end{exercise}

\begin{solution}
	\begin{proof}
		We prove this by first independently proving the forward and reverse directions of this statement.

		\begin{steps}
			\item Prove the forward direction ($T_1$ is denumerable $\implies$ there exists a bijection from $T_1$ to a denumerable set $T_2$).
		\end{steps}

		Since both $T_1$ and $T_2$ are denumerable, we know that there exists two bijection functions $f$, $g$ such that $f: \mathbb{N} \rightarrow T_1$ and $g: \mathbb{N} \rightarrow T_2$ respectively. However, we know that $g^{-1}$ is also a bijective function from $T_2$ onto $\mathbb{N}$, and from exercise 1.1.21, we know that the composite function $g^{-1} \circ f$ is a bijection from $T_1$ onto $T_2$. Thus, if $T_1$ and $T_2$ are both denumerable, there must exist a bijection from $T_1$ to $T_2$.

		\begin{steps}
			\item Prove the reverse direction (A bijection from $T_1$ to a denumerable set $T_2 \implies T_1$ is denumerable),
		\end{steps}

		We know that there exists a bijection from $T_1$ to $T_2$. Let's call such a bijection, $f: T_1 \rightarrow T_2$. Since $T_2$ is denumerable, we know that there exists a bijection from $T_2$ onto $\mathbb{N}$. Let's call that bijection $g: T_2 \rightarrow \mathbb{N}$. However, from exercise 1.1.21, we proved that $g \circ f$ would be a bijective map from $T_1$ onto $\mathbb{N}$, making $(f \circ g)^{-1}$ a bijection from $\mathbb{N}$ onto $T_1$. Thus, by definition of denumerability, $T_1$ is denumerable. \\\\

		Thus, since both the forward and reverse directions are true, $T_1$ is denumerable $\iff$ there exists a bijection from $T_1$ to a denumerable set $T_2$.
	\end{proof}
\end{solution}

\begin{exercise}[13]
	Prove that the collection $\mathcal{F}(\mathbb{N})$ of all \emph{finite} subsets of $\mathbb{N}$ is countable.
\end{exercise}

\begin{solution}
	\begin{proof}
		\underline{Note}: this proof is technically incorrect since certain numbers, such as 4, map to the singleton set $\{2\}$, which is undesirable, but I feel like this approach is in the right direction. Will update with a more correct proof in future.

		To prove that $\mathcal{F}$ is bijective, consider the fact that every natural number (besides 1) can be decomposed into the product of primes. Using this fact, we can construct a bijection between $\mathbb{N}$ and $\mathcal{F}$ with the function:

		\begin{align*}
			f(n) &=
			\begin{cases}
				\{1\} \; \text{if} \; n = 1 \\
				\{i + 1\} \; \text{if} \; n \; \text{is the} \; i^{th} \; \text{prime number.} \\
				\{j, k\} \; \text{if} \; n \; \text{has two prime factors} \; n_1, n_2 \; \text{and where}\\ \; n_1, n_2 \; \text{are the} \; j^{th}\ \text{and} \; k^{th} \; \text{prime numbers respectively.} \\
				\vdots
			\end{cases}
		\end{align*}

		Thus, since there exists a bijective mapping from $\mathbb{N}$ to $\mathcal{F}$, by definition of denumerability, $\mathcal{F}$ is denumerable and is thus countable.

	\end{proof}
\end{solution}

\ifdefined\iscombined
\newpage
\else
  \end{document}
\fi
